\input{../tex-inputs/latex-leseansicht-vorspann}

\section[Gustav Schwarzkopf an Arthur Schnitzler, 26. 4. 1899]{L04301 Gustav Schwarzkopf an Arthur Schnitzler, 26. 4. 1899}
\nopagebreak\mylabel{L04301v}
\rehead{ }\normalsize\beginnumbering\briefempfaengerindex{Schnitzler, Arthur@\textsc{Schnitzler, Arthur}!zzzSchwarzkopf, Gustav@\emph{von Gustav Schwarzkopf}!1899-04-262@{26. 4. 1899}|(be}
\toendnotes[C]{\smallbreak\pagebreak[2]}
\correspDesc{Versand  durch Gustav Schwarzkopf am 26. 4. 1899 in Wien
\newline{}Übermittlung  am 26. 4. 1899 in Wien
\newline{}Erhalt  durch Arthur Schnitzler im Zeitraum [27. 4. 1899 – 30. 4. 1899?] in Berlin}\toendnotes[C]{\smallbreak}
\Standort{CUL, Schnitzler, B 96a.}
\physDesc{Postkarte, 605 Zeichen
\newline{}Handschrift: schwarze Tinte, deutsche Kurrent
\newline{}Versand: Stempel: »\nobreak{}\oindex{I., Innere Stadt@\textbf{I., Innere Stadt}, \emph{Verwaltungsgebiet}|pwk}Wien 1/1 1, 26. 4. 99, 7–8N\nobreak{}«.  }\toendnotes[C]{\smallbreak}\pstart{}{\pb}Herrn \textsc{D\textsuperscript{r} Arthur Schmitzler}\pend{}\pstart{}\textsc{Berlin\oindex{Berlin@\textbf{Berlin}, \emph{Hauptstadt}|pw}}\pend{}\pstart{}\textsc{Savoy-Hôtel\oindex{Hotel Savoy [Berlin]@\textbf{Hotel Savoy [Berlin]}, \emph{Hotel}|pw}}\pend{}{\bigskip}\vspace{1em}
\pstart
           \raggedleft{}{\pb}Wien\oindex{Wien@\textbf{Wien}, \emph{Verwaltungsgebiet}|pw}{ }26. IV. 99\pend
           \vspace{0.5em}
\pstart
           Lieber Arthur, ich will Ihnen nur für Ihre \label{K_L04301-1v}\edtext{Karte}{\lemma{\textnormal{\emph{Karte}}}\Cendnote{\textnormal{{XXXX ref} XXXX
               }}}\label{K_L04301-1} danken und Ihre
               Grüße herzlichſt erwidern.\pend
           
\pstart
           Ereignet hat sich nichts. Das Wetter iſt noch i{\geminationm}er{ }ſo{ }ſchlecht, als ob wir die Abſicht hätten, eine größere Partie mit der Stadtbahn\orgindex{Wiener Stadtbahn@Wiener Stadtbahn|pw} zu machen und die \label{K_L04301-2v}\edtext{dritte Nu{\geminationm}er
               der »Fackel\pwindex{Fackel@\emph{Die Fackel}|pw}«}{\lemma{\textnormal{\emph{dritte … »Fackel«}}}\Cendnote{\textnormal{Am Titelblatt ist die dritte Nummer mit »Ende April« datiert. Spätestens am 4. 5. 1899 wurde
               das Heft ausgegeben.}}}\label{K_L04301-2}, auf welche alle Graben\oindex{Wien@\textbf{Wien}!I., Innere Stadt@\textbf{I., Innere Stadt}!Graben@\textbf{Graben}, \emph{Straße}|pw}ſpaziergänger mit Ungeduld warten, iſt noch nicht
               erſchienen. We{\geminationn} ich Ihnen{ }ſage, daß ich mit meinen
               Nachmittagen nichts anzufangen weiß und daß ich mich{ }ſehr gerne mit Ihnen zuſa{\geminationm}en in das »Berliner\oindex{Berlin@\textbf{Berlin}, \emph{Hauptstadt}|pw}
               Straßenleben“{ }ſtürzen möchte,{ }ſo dürfen Sie das glauben\pend
           
\pstart
           Ihrem{\\[\baselineskip]}\spacefill\mbox{Gustav.}\pend
           \leftskip=0em{}\selectlanguage{ngerman}\endnumbering\briefempfaengerindex{Schnitzler, Arthur@\textsc{Schnitzler, Arthur}!zzzSchwarzkopf, Gustav@\emph{von Gustav Schwarzkopf}!1899-04-262@{26. 4. 1899}|)be}\mylabel{L04301h}
\begin{anhang}
\end{anhang}\newcommand{\dateiname}{L04301}\newcommand{\titel}{Gustav Schwarzkopf an Arthur Schnitzler, 26. 4. 1899}\newcommand{\editorInnen}{Herausgegeben von Jahnke, SelmaMüller, Martin Anton}\input{../tex-inputs/latex-leseansicht-abspann}
